% ch4.tex -- Module 7 ESSENTIALS, Chapter 4: Light and special relativity
% Condensed from src/mod7/ch4.tex (final full book). Knowledge points only:
% no worked examples, quick checks, exam tips or questions. All diagrams are TikZ copied
% (some shrunk) from the full chapter, so no figure files are needed.
\chapter{Light and Special Relativity}\label{ch4:chap}
\covers{M7.IQ4.01, M7.IQ4.02, M7.IQ4.03, M7.IQ4.04}

\begin{summarybox}{In this chapter}
\begin{itemize}
  \item Einstein's two postulates and the evidence for and against them.
  \item Simultaneity, time dilation and length contraction; muons in two frames.
  \item Relativistic momentum, and why nothing with mass can reach $c$.
  \item $E = mc^2$ for the Sun, annihilation and burning fuel; the key experiments.
\end{itemize}
\end{summarybox}

%% ==================================================================
\section{Frames of reference and Einstein's postulates}\label{ch4:sec:postulates}
\covers{M7.IQ4.01}

\begin{bigidea}
Motion is always measured from a viewpoint. Einstein said that no steady viewpoint is special,
and that everyone measures the same speed for light in a vacuum.
\end{bigidea}

\begin{keywords}
  \kw{Frame of reference}{a viewpoint, with its own rulers and clocks.}
  \kw{Inertial frame}{a frame that is not accelerating: at rest or at constant velocity.}
  \kw{Event}{something that happens at one place and one instant, such as a flash.}
  \kw{Postulate}{a starting statement that a theory takes as true.}
\end{keywords}

\begin{itemize}
  \item Inertial means constant \emph{velocity}. A \term{non-inertial frame} speeds up, slows
        down or turns. The Earth's surface is usually treated as inertial.
  \item \term{Galilean relativity}: the laws of \emph{motion} are the same in every inertial
        frame, and everyday velocities add. Einstein's postulate 1 extends it to all of physics.
  \item Maxwell's equations give one speed for light, $c = \qty{3.00e8}{m.s^{-1}}$, but not
        \emph{relative to what}. Light was thought to need a \term{medium} (a substance that
        carries a wave): an invisible \term{luminiferous aether}.
  \item The Earth moves through space at about \qty{30}{km.s^{-1}}, so light's speed should
        change with direction. Michelson and Morley (1887) found no change: a \term{null
        result} (no effect found).
\end{itemize}

\subsection{Einstein's two postulates (1905)}
\term{Special relativity} is the theory of how space and time behave for inertial observers.
It rests on two postulates. Learn their exact wording.
\begin{enumerate}
  \item \textbf{All inertial frames of reference are equivalent.} The laws of physics, including
        those of light, are the same in every inertial frame. No experiment can show which
        inertial frame is ``really'' at rest.
  \item \textbf{The speed of light in a vacuum is an absolute constant.} Every inertial observer
        measures $c = \qty{3.00e8}{m.s^{-1}}$, however fast the source or the observer moves.
\end{enumerate}
\begin{itemize}
  \item By postulate 1, Maxwell's equations give the same $c$ in every inertial frame.
  \item Velocities do \emph{not} add for light: a headlight beam travels at $c$, not $c + v$. So
        observers moving relative to each other must disagree about distance, time or both.
  \item $c$ is constant only \emph{in a vacuum}. Light slows down in glass and water.
\end{itemize}

\begin{table}
  \caption{Evidence about Einstein's postulates.}\label{ch4:tab:evidence}
  \centering\small
  \begin{tabular}{@{}L{42mm}L{78mm}L{18mm}@{}}
    \toprule
    \colhead{Evidence} & \colhead{What was found} & \colhead{Postulate}\\
    \midrule
    Michelson--Morley (1887) & null result: the same speed in every direction & 1 and 2\\
    No absolute motion found & the same laws of physics on the moving Earth & 1\\
    Fast pions (short-lived particles), CERN (1964) & light from a source at $0.99975c$
      still travels at $c$ & 2\\
    Modern laser tests & $c$ is the same in all directions to 1 part in $10^{17}$ & 2\\
    Muons, atomic clocks, GPS & every prediction of relativity confirmed & 1 and 2\\
    OPERA (2011) claim & ``faster-than-light'' neutrinos (tiny, almost massless particles):
      really a loose cable & neither\\
    \bottomrule
  \end{tabular}
\end{table}

\begin{itemize}
  \item \textbf{Judgement} (\cref{ch4:tab:evidence}): strongly supported. Many tests agree; no
        result against them has survived checking. (Evidence supports a theory; it does not
        prove it.)
  \item \textbf{Limitation:} Michelson and Morley's source was at rest. The pion test covers a
        moving source.
\end{itemize}

%% ==================================================================
\section{Time dilation and length contraction}\label{ch4:sec:time}
\covers{M7.IQ4.02}

\begin{bigidea}
If everyone measures the same speed for light, moving observers must disagree about time and
length. A moving clock runs slow. A moving object is shorter along its motion.
\end{bigidea}

\begin{keywords}
  \kw{Proper time ($t_0$)}{time between two events, measured by a clock present at both.}
  \kw{Proper length ($l_0$)}{length measured in the frame where the object is at rest.}
\end{keywords}

A \term{thought experiment} is an imagined experiment, worked through with logic. It explores
speeds near $c$ that nobody could test in 1905. It is \emph{not} evidence: it shows what the
theory predicts, and real experiments test it.

\subsection{Thought experiment 1: lightning and a fast train}
\begin{sidebyside}[0.5]
  \centering
  \begin{tikzpicture}[x=0.62cm, y=0.62cm]
    \foreach \yy/\xx/\lab in {0/0/{(1) Lightning strikes both ends.},
        -2.7/1.333/{(2) The front flash reaches Peter first.}} {
      \node[diagramlabel, font=\scriptsize, anchor=west] at (-0.7,\yy+1.6) {\lab};
      \draw[hscink, line width=0.8pt] (-0.7,\yy) -- (10.2,\yy);
      \foreach \mx/\ml in {0/A, 8/B} {
        \draw[hscmuted, line width=1pt] (\mx-0.12,\yy-0.12) -- (\mx+0.12,\yy+0.12)
          (\mx-0.12,\yy+0.12) -- (\mx+0.12,\yy-0.12);
        \node[diagramlabel, font=\scriptsize, text=hscmuted] at (\mx,\yy-0.36) {\ml};
      }
      \draw[object] (\xx,\yy+0.12) rectangle (\xx+8,\yy+0.9);
      \node[diagramlabel, font=\scriptsize, anchor=west, inner sep=1pt] at (\xx+0.05,\yy+0.51) {rear};
      \node[diagramlabel, font=\scriptsize, anchor=east] at (\xx+7.88,\yy+0.51) {front};
      \fill[hscvel] (\xx+4,\yy+0.42) circle (0.09);
      \node[diagramlabel, font=\scriptsize, text=hscvel, anchor=east] at (\xx+3.85,\yy+0.45) {Peter};
      \fill[hscforce] (4,\yy-0.42) circle (0.09);
    }
    \foreach \bx in {0,8} {
      \draw[hscacc, line width=1.1pt] (\bx+0.15,1.35) -- (\bx-0.05,1.12) -- (\bx+0.12,1.09) -- (\bx-0.02,0.92);
    }
    \draw[velocity] (8.3,0.5) -- (9.3,0.5) node[right, text=hscvel] {$v$};
    \node[diagramlabel, font=\scriptsize, text=hscforce, anchor=west] at (4.1,-0.48) {Olivia};
    \node[diagramlabel, font=\scriptsize, text=hscforce, anchor=north, inner sep=1pt] at (4,-3.27) {Olivia};
    \draw[hscray, line width=1.4pt] (2.667,-3.3) -- (2.667,-1.7);
    \draw[hscray, line width=1pt, ->] (2.667,-1.65) -- ++(0.55,0);
    \draw[hscray, line width=1.4pt] (5.333,-3.3) -- (5.333,-1.7);
    \draw[hscray, line width=1pt, ->] (5.333,-1.65) -- ++(-0.55,0);
  \end{tikzpicture}
  \captionof{figure}{Lightning and a fast train, in the platform's frame. Orange: the flash
    fronts.}\label{ch4:fig:lightning}
\nextside
  \begin{itemize}
    \item Olivia (on the platform, halfway between A and B): both flashes reach her together
          (\cref{ch4:fig:lightning}). For her the strikes are \term{simultaneous} (at the same
          time).
    \item The train carries Peter towards the front flash, so it reaches him first.
    \item Peter is at rest mid-train and both flashes travel at $c$. So, for him, the front
          strike \emph{really happened first}. His frame is just as valid (postulate 1).
  \end{itemize}
\end{sidebyside}
This is the \term{relativity of simultaneity}: events at different places that are simultaneous
in one frame are not simultaneous in a frame moving along the line between them.

\subsection{Thought experiment 2: the light clock}
\begin{sidebyside}[0.5]
  \centering
  \begin{tikzpicture}[x=0.82cm, y=0.82cm]
    \draw[hscink, line width=2pt] (-0.5,0) -- (0.5,0) (-0.5,2) -- (0.5,2);
    \node[diagramlabel, font=\scriptsize, anchor=south] at (0,2.08) {mirror};
    \node[diagramlabel, font=\scriptsize, anchor=north] at (0,-0.08) {mirror};
    \draw[ray, ->] (-0.13,0.06) -- (-0.13,1.94);
    \draw[ray, ->] (0.13,1.94) -- (0.13,0.06);
    \draw[<->, hscmuted, line width=0.6pt] (0.8,0) -- node[right, text=hscink] {$d$} (0.8,2);
    \node[diagramlabel, font=\footnotesize] at (0.1,-1.0) {(a) at rest};
    \begin{scope}[xshift=2.0cm]
      \foreach \x/\op in {0/25, 2.2/50, 4.4/100}
        \draw[hscink!\op, line width=2pt] (\x-0.5,0) -- (\x+0.5,0) (\x-0.5,2) -- (\x+0.5,2);
      \draw[ray] (0,0.04) -- (2.2,1.96);
      \draw[ray] (2.2,1.96) -- (4.4,0.04);
      \draw[guide] (0,-0.08) -- (2.2,-0.08) -- (2.2,1.96);
      \node[diagramlabel, font=\footnotesize, below] at (1.1,-0.12) {$vt/2$};
      \node[diagramlabel, font=\footnotesize, right] at (2.22,1.1) {$d$};
      \node[diagramlabel, font=\footnotesize, above left, text=hscray] at (1.0,1.0) {$ct/2$};
      \draw[velocity] (0.6,2.45) -- (1.8,2.45) node[right, text=hscvel] {$v$};
      \node[diagramlabel, font=\footnotesize] at (2.2,-1.0) {(b) moving at speed $v$};
    \end{scope}
  \end{tikzpicture}
  \captionof{figure}{A light clock in its own frame (a) and in a frame where it moves
    (b).}\label{ch4:fig:lightclock}
\nextside
  \begin{itemize}
    \item A \term{light clock}: light bouncing between two mirrors (\cref{ch4:fig:lightclock}).
          (a) At rest, one tick takes $t_0 = 2d/c$.
    \item (b) Moving: the slanted path is longer, but the light still travels at $c$. So each
          tick takes longer: \term{time dilation}.
    \item Pythagoras for the triangle in (b): $(ct/2)^2 = d^2 + (vt/2)^2$. With
          $t_0 = 2d/c$, this gives $t = t_0/\sqrt{1 - v^2/c^2}$.
  \end{itemize}
\end{sidebyside}

\subsection{Proper time and proper length}
\begin{itemize}
  \item $t_0$ belongs to the clock present at \emph{both} events, in whatever frame: the muon's
        own clock for its life, the ship's clock for a trip. Others measure a longer $t$.
  \item \term{Length contraction}: a moving object is measured shorter along its motion, not
        across it.
        For the Earth-to-star distance, $l_0$ is measured in the Earth--star frame.
  \item Both effects work both ways. Your own clock and ruler always seem normal to you. The
        effects are tiny at everyday speeds and large near $c$.
\end{itemize}

\subsection{Muons: both effects in action}
\begin{sidebyside}[0.46]
  \centering
  \begin{tikzpicture}[x=0.72cm, y=0.72cm]
    \fill[ground] (0,-0.25) rectangle (2.6,0);
    \draw[hscink, line width=0.8pt] (0,0) -- (2.6,0);
    \draw[hscmuted, dashed, line width=0.6pt] (0,3.6) -- (2.6,3.6);
    \draw[<->, hscmuted, line width=0.6pt] (-0.2,0) -- node[left, text=hscink, font=\footnotesize]
      {$l_0$} (-0.2,3.6);
    \node[diagramlabel, font=\scriptsize, anchor=west] at (2.65,-0.12) {ground};
    \node[diagramlabel, font=\scriptsize, anchor=south, align=center] at (1.3,3.62)
      {top of\\[-2pt]atmosphere};
    \fill[object] (1.0,3.3) circle (0.11);
    \node[diagramlabel, font=\scriptsize, anchor=west] at (1.15,3.3) {muon};
    \draw[velocity] (1.0,3.05) -- (1.0,1.9) node[left, text=hscvel, font=\scriptsize] {$v$};
    \node[diagramlabel, font=\scriptsize, align=left, anchor=west] at (1.15,2.4)
      {its clock\\[-2pt]runs slow};
    \node[diagramlabel, font=\footnotesize] at (1.3,-0.7) {(a) Earth frame};
    \begin{scope}[xshift=2.9cm]
      \fill[ground] (0,1.95) rectangle (2.6,2.2);
      \draw[hscink, line width=0.8pt] (0,2.2) -- (2.6,2.2);
      \node[diagramlabel, font=\scriptsize, anchor=west] at (2.65,2.08) {ground};
      \draw[hscmuted, dashed, line width=0.6pt] (0,3.2) -- (2.6,3.2);
      \node[diagramlabel, font=\scriptsize, anchor=south, align=center] at (1.3,3.22)
        {top of\\[-2pt]atmosphere};
      \draw[<->, hscmuted, line width=0.6pt] (-0.2,2.2) -- node[left, text=hscink, font=\footnotesize]
        {$l$} (-0.2,3.2);
      \fill[object] (1.0,2.95) circle (0.11);
      \node[diagramlabel, font=\scriptsize, anchor=west] at (1.15,2.92) {muon at rest};
      \draw[velocity] (1.8,0.55) -- (1.8,1.75);
      \node[diagramlabel, font=\scriptsize, align=right, anchor=east] at (1.65,1.1) {Earth\\[-2pt]rushes up};
      \node[diagramlabel, font=\footnotesize] at (1.3,-0.7) {(b) muon frame};
    \end{scope}
  \end{tikzpicture}
  \captionof{figure}{One muon trip in two frames (not to scale).}\label{ch4:fig:muons}
\nextside
  \begin{itemize}
    \item Cosmic rays (fast protons from space) make \term{muons} (like heavy electrons)
          10--15 km up. Muons are \term{unstable} (they decay):
          mean lifetime at rest \qty{2.2}{\micro\second}, only about \qty{660}{m} of travel
          ($s = vt$) without relativity. Yet many reach the ground (\cref{ch4:fig:muons}).
    \item (a) Earth frame: the muon's clock runs slow, so it lives long enough to cross $l_0$.
    \item (b) Muon frame: it lives $t_0$, but the atmosphere is contracted to $l$.
    \item Both frames agree. Never mix dilated time and contracted length in one frame.
  \end{itemize}
\end{sidebyside}
\begin{itemize}
  \item \textbf{Evidence:} Frisch and Smith (1963) counted muons near $0.995c$: 563 per hour on
        a mountain \qty{1.9}{km} up, 408 per hour at sea level. Without time dilation only about
        30 would survive. The counts give a dilation factor of about 9: the factor is 10 at
        $0.995c$, but the muons slow down in the air.
  \item \textbf{Evidence:} at CERN (1977), muons stored in a ring at $0.999418c$ lived 29.3
        times longer than at rest, as time dilation predicts.
\end{itemize}

% layout (final page check): keep the formulas with their symbol table on one page
\begin{formulacard}[unbreakable]{Time dilation and length contraction}
  \cardformula{t = \frac{t_0}{\sqrt{\left(1 - \dfrac{v^2}{c^2}\right)}} \qquad\qquad
    l = l_0\sqrt{\left(1 - \dfrac{v^2}{c^2}\right)}}
  \begin{symbols}
    \sym{$t$}{time measured in a frame where the clock moves}{\si{s}}
    \sym{$t_0$}{proper time: by a clock present at both events}{\si{s}}
    \sym{$l$}{length measured in a frame where the object moves}{\si{m}}
    \sym{$l_0$}{proper length: measured where the object is at rest}{\si{m}}
    \sym{$v$}{relative speed of the clock or object and the observer}{\si{m.s^{-1}}}
    \sym{$c$}{speed of light in a vacuum, \qty{3.00e8}{m.s^{-1}}}{\si{m.s^{-1}}}
  \end{symbols}
  \whentouse{$t$: the same two events seen from two inertial frames. $l$: lengths \emph{along}
    the motion.}
  \trap{square $\dfrac{v}{c}$ first: $\sqrt{1 - 0.9^2}$, not $\sqrt{1 - 0.9}$. Always
    $t > t_0$ and $l < l_0$.}
\end{formulacard}

%% ==================================================================
\section{Relativistic momentum and the speed limit}\label{ch4:sec:momentum}
\covers{M7.IQ4.03}

\begin{bigidea}
Near $c$, momentum grows much faster than $m_0v$ and heads towards infinity. So nothing with
mass can ever reach the speed of light.
\end{bigidea}

\begin{sidebyside}[0.45]
  \centering
  \begin{tikzpicture}
    \begin{axis}[hsc, width=\linewidth, height=0.72\linewidth,
        xmin=0, xmax=1.08, ymin=0, ymax=7, xtick={0,0.2,...,1.0},
        xticklabel style={/pgf/number format/.cd, fixed, precision=1},
        xlabel={speed $v/c$}, ylabel={momentum $p/(m_0c)$}]
      \addplot[hscseriesA, no marks, line width=1.2pt, domain=0:0.99, samples=200] {x/sqrt(1-x^2)};
      \addplot[hscseriesB, no marks, dashed, line width=1.2pt, domain=0:1.0, samples=2] {x};
      \draw[guide] (axis cs:1,0) -- (axis cs:1,7);
      \node[font=\footnotesize, anchor=east, text=hscseriesA] at (axis cs:0.92,5.2) {$p_\mathrm{v}$};
      \node[font=\footnotesize, anchor=north, inner sep=0pt, text=hscseriesB] at (axis cs:0.93,0.67) {$m_0v$};
      \node[font=\footnotesize, anchor=north, rotate=90] at (axis cs:1.02,3.5) {$v = c$};
    \end{axis}
  \end{tikzpicture}
  \captionof{figure}{Momentum against speed. The dashed line is
    $p = m_0v$.}\label{ch4:fig:momentum}
\nextside
  \begin{itemize}
    \item Near $c$, $p = mv$ fails: momentum would be conserved in one inertial frame but not in
          another (against postulate 1). The \term{relativistic momentum}
          $p_\mathrm{v}$ uses the \term{rest mass} $m_0$ (mass measured at rest).
    \item As $v \to c$ (``gets closer to $c$''), the root $\to 0$, so $p_\mathrm{v} \to \infty$
          (\cref{ch4:fig:momentum}).
    \item A finite force for a finite time gives only a finite $\Delta p = F_\mathrm{net}\Delta t$.
          So $v$ never reaches $c$.
  \end{itemize}
\end{sidebyside}

\begin{itemize}
  \item Below about $0.1c$, $p_\mathrm{v} \approx m_0v$. A photon has no rest mass and always
        travels at $c$.
  \item \textbf{Evidence:} Bertozzi's (1964) electrons of up to \qty{15}{\mega\eV} approached $c$
        but never passed it. The Large Hadron Collider (LHC, CERN's biggest accelerator) gives
        protons 7000 times their rest energy, but they stay below $c$.
  \item A magnet bends a moving charge into a circle: \hscbadgenotonsheet\ $r = p/(qB)$
        (from $qvB = mv^2/r$).
  \item In a \term{synchrotron} (a ring-shaped accelerator) $r$ is fixed, so $B$ must rise
        with $p_\mathrm{v}$, not $m_0v$: a daily test of relativity.
  \item The HSC formula keeps $m_0$ fixed: it is the momentum that grows, not the mass.
\end{itemize}

\begin{formulacard}{Relativistic momentum}
  \cardformula{p_\mathrm{v} = \frac{m_0 v}{\sqrt{\left(1 - \dfrac{v^2}{c^2}\right)}}}
  \begin{symbols}
    \sym{$p_\mathrm{v}$}{momentum of an object moving at speed $v$}{\si{kg.m.s^{-1}}}
    \sym{$m_0$}{rest mass}{\si{kg}}
    \sym{$v$}{speed of the object relative to the observer}{\si{m.s^{-1}}}
    \sym{$c$}{speed of light in a vacuum, \qty{3.00e8}{m.s^{-1}}}{\si{m.s^{-1}}}
  \end{symbols}
  \whentouse{any moving object with mass. Not for a photon (no rest mass).}
  \trap{on the top line, $v$ is in \si{m.s^{-1}} (fraction of $c$ times \num{3.00e8}).}
\end{formulacard}

%% ==================================================================
\section{Mass--energy equivalence: \texorpdfstring{$E = mc^2$}{E = mc2}}\label{ch4:sec:emc2}
\covers{M7.IQ4.04}

\begin{bigidea}
Mass is a very concentrated form of energy. When energy is released, a little mass is lost:
the energy equals the lost mass times $c^2$.
\end{bigidea}

\begin{keywords}
  \kw{Mass defect ($\Delta m$)}{mass before minus mass after: the mass that disappears.}
  \kw{Antiparticle}{a particle with the same mass as another but the opposite charge.}
  \kw{Annihilation}{a particle and its antiparticle meet and become photons.}
\end{keywords}

\begin{itemize}
  \item Any mass $m$ has a \term{rest energy} $E = mc^2$. Since
        $c^2 = \qty{9.00e16}{m^2.s^{-2}}$ is huge, a tiny mass is equivalent to a huge energy
        (one gram: about a day's output of a large power station).
  \item Every process that releases energy loses mass. A sealed box gets lighter only as the
        energy escapes. In reverse, collision energy can make new particles.
  \item Mass in atomic mass units (\si{\amu}): multiply by 931.5 to get MeV. Mass in kg:
        $E = mc^2$ gives J ($\qty{1}{\eV} = \qty{1.602e-19}{J}$). Do not mix the two routes.
  \item \textbf{The Sun:} \term{nuclear fusion} (light nuclei join) turns four hydrogen nuclei
        into helium-4. About 0.7\,\% of the mass becomes energy; the Sun loses 4 million tonnes a
        second.
  \item \textbf{Burning fuel:} the mass lost is only about 5 parts in 10 billion of the fuel's
        mass: far too small to weigh. We know it only by calculation.
\end{itemize}

\begin{sidebyside}[0.5]
  \centering
  \begin{tikzpicture}[x=0.72cm, y=0.7cm]
    \fill[hscvel!80] (0,0) circle (0.2) node[white, font=\footnotesize\bfseries] {$-$};
    \fill[hscforce!85] (2.4,0) circle (0.2) node[white, font=\footnotesize\bfseries] {$+$};
    \draw[velocity] (0.3,0) -- (0.95,0);
    \draw[velocity] (2.1,0) -- (1.45,0);
    \node[diagramlabel, font=\scriptsize, below=2.5mm] at (0,0) {electron};
    \node[diagramlabel, font=\scriptsize, below=2.5mm] at (2.4,0) {positron};
    \draw[->, hscmuted, line width=1pt] (3.1,0) -- (3.8,0);
    \begin{scope}[xshift=4.5cm]
      \fill[hscacc] (0,0) circle (0.07);
      \foreach \dir in {1,-1}
        \draw[hscray, line width=1.1pt, decorate, decoration={snake, amplitude=1.1mm,
          segment length=3mm, post length=2.5mm}, ->] (\dir*0.15,0) -- (\dir*2.0,0);
      \node[diagramlabel, font=\scriptsize, above=2mm, text=hscray] at (1.2,0) {photon};
      \node[diagramlabel, font=\scriptsize, above=2mm, text=hscray] at (-1.2,0) {photon};
    \end{scope}
  \end{tikzpicture}
  \captionof{figure}{Annihilation (pair almost at rest): total momentum is zero before and
    after.}\label{ch4:fig:annihilation}
\nextside
  \textbf{Annihilation:} an electron meets a \term{positron} ($\mathrm{e}^+$, its antiparticle).
  \emph{All} their mass ($2m_\mathrm{e}$) becomes two equal gamma-ray photons in opposite
  directions (\cref{ch4:fig:annihilation}). Hospital PET scanners (a type of medical scan)
  detect these pairs.
\end{sidebyside}

\textbf{Energy released in annihilation} \hscbadgecombine{$E = mc^2$}\
$E = 2m_\mathrm{e}c^2$. So each of the two photons has $E = m_\mathrm{e}c^2$.

\begin{formulacard}{Mass--energy equivalence}
  \cardformula{E = mc^2}
  \begin{symbols}
    \sym{$E$}{energy equivalent to the mass (energy released)}{\si{J}}
    \sym{$m$}{mass; use $\Delta m$, the mass lost, for energy released}{\si{kg}}
    \sym{$c$}{speed of light in a vacuum, \qty{3.00e8}{m.s^{-1}}}{\si{m.s^{-1}}}
  \end{symbols}
  \whentouse{energy released when mass is lost (fusion, annihilation, burning). For a star or a
    power station of power output $P$ (in W): $\dfrac{\Delta m}{\Delta t} = \dfrac{P}{c^2}$
    \hscbadgecombine{$P = \frac{\Delta E}{\Delta t}$ and $E = mc^2$}}
  \trap{square $c$. Use only the mass that is \emph{lost} (in annihilation, all of it).}
\end{formulacard}

%% ==================================================================
% layout: keep this heading with the (unbreakable) Michelson--Morley box below it
\Needspace{26\baselineskip}
\section{Key experiments}\label{ch4:sec:experiments}
\covers{M7.IQ4.01, M7.IQ4.02}

\begin{experiment}[unbreakable]{Michelson and Morley (1887): looking for the aether}
\aim{To detect the Earth's motion through the aether.}
\method{
  \begin{sidebyside}[0.4]
    \centering
    \begin{tikzpicture}[x=0.5cm, y=0.5cm]
      \fill[hscmuted!12, rounded corners=2mm] (-0.9,-2.9) rectangle (6.9,3.5);
      \draw[object] (-0.7,-0.3) rectangle (0,0.3);
      \node[diagramlabel, font=\scriptsize, above] at (-0.35,0.3) {source};
      \draw[ray] (0,0.06) -- (2.95,0.06);
      \draw[hscink, line width=1.6pt] (2.65,-0.35) -- (3.35,0.35);
      \node[diagramlabel, font=\scriptsize, anchor=north east, align=right] at (2.8,-0.35)
        {half-silvered\\[-2pt] mirror};
      \draw[ray] (3.05,0.06) -- (5.9,0.06);
      \draw[ray] (5.9,-0.08) -- (3.08,-0.08);
      \draw[hscink, line width=2pt] (6.0,-0.5) -- (6.0,0.5);
      \node[diagramlabel, font=\scriptsize, anchor=south] at (6.0,0.55) {mirror};
      \draw[ray] (2.94,0.1) -- (2.94,2.9);
      \draw[ray] (3.08,2.9) -- (3.08,0.1);
      \draw[hscink, line width=2pt] (2.5,3.0) -- (3.5,3.0);
      \node[diagramlabel, font=\scriptsize, anchor=west] at (3.55,3.0) {mirror};
      \draw[ray] (3.0,-0.12) -- (3.0,-1.75);
      \draw[object] (2.65,-2.35) rectangle (3.35,-1.8);
      \node[diagramlabel, font=\scriptsize, anchor=west] at (3.4,-2.05) {telescope};
    \end{tikzpicture}
    \captionof{figure}{The interferometer, from above.}\label{ch4:fig:mm}
  \nextside
    An \term{interferometer} (\cref{ch4:fig:mm}) split one beam into two at right angles. Each
    went along an arm and back, and they recombined to make interference fringes (bright and
    dark bands). The apparatus floated on mercury and was slowly rotated, at different times of
    day and year. A change in light's speed would shift the fringes.
  \end{sidebyside}}
\results{Predicted shift: about 0.4 of a fringe. Seen: less than 0.02 of a fringe (null result).}
\conclusion{No aether wind: light has the same speed in all directions on the moving Earth.}
\significance{No evidence for an aether. The result agrees with both postulates.}
\end{experiment}

\begin{experiment}{Hafele--Keating atomic clocks (1971)}
\aim{To test time dilation with very accurate atomic clocks flown around the world.}
\method{Four caesium \term{atomic clocks} flew around the world eastward, then westward, and
  were compared with ground clocks. Seen from the Earth's centre, the eastward clock moves fastest
  (the Earth spins east), so it runs slowest.}
\results{Eastward clocks lost \qty{59}{ns} (predicted \qty{40}{ns}); westward clocks gained
  \qty{273}{ns} (predicted \qty{275}{ns}). The predictions include a gravity effect (general
  relativity, not in the HSC course).}
\conclusion{Measured and predicted times agree within the uncertainties: moving clocks run slow,
  even at aeroplane speeds.}
\significance{Direct evidence for time dilation with real clocks at everyday speeds. GPS
  satellites need similar corrections: speed and gravity together, about
  \qty{38}{\micro\second} a day.}
\end{experiment}

\begin{experiment}{Evidence from cosmology: slow-motion supernovae}
\aim{To test whether events in very distant galaxies are measured to run slow.}
\method{Every \term{Type Ia supernova} (exploding star) fades in nearly the same way. Hundreds
  were timed, and each \term{redshift} $z$ (stretch of the light) was measured:
  \hscbadgenotonsheet\ $1 + z = \dfrac{\lambda_\mathrm{observed}}{\lambda_\mathrm{emitted}}$.}
\results{Distant ones fade more slowly, by the factor $1 + z$ (twice as long at $z = 1$).}
\conclusion{\term{Cosmological time dilation}, mainly from the expansion of the Universe
  (general relativity).}
\significance{Evidence for time dilation from cosmological studies: time is not absolute, even
  across the Universe.}
\end{experiment}

\begin{summarybox}{Remember}
\begin{itemize}
  \item Postulate 1: all inertial frames of reference are equivalent. Postulate 2: the speed of
        light in a vacuum is an absolute constant.
  \item Simultaneity is relative. Moving clocks run slow ($t > t_0$); moving objects are shorter
        along the motion ($l < l_0$). Muons show both.
  \item Proper time $t_0$: by a clock present at both events. Proper length $l_0$: measured where
        the object is at rest.
  \item As $v \to c$, $p_\mathrm{v} \to \infty$, so nothing with mass reaches $c$.
  \item $\Delta E = \Delta m\,c^2$: the Sun (0.7\,\%), annihilation (all the mass), fuel (tiny).
  \item Evidence: Michelson--Morley, muons, atomic clocks, accelerators, supernovae.
\end{itemize}
\end{summarybox}
